\documentclass[12pt]{report}
%\usepackage{helvet} 
\usepackage[a4paper,margin=1in]{geometry}
\usepackage[T1]{fontenc}
\usepackage{tgheros} 
\usepackage[utf8]{inputenc}
\usepackage{lmodern}
\usepackage{amsmath,amssymb,amsfonts,mathtools}
\usepackage{bm}
\usepackage{microtype}
\usepackage{xcolor}
\usepackage{enumitem}
\usepackage{booktabs}
\usepackage{longtable}
\usepackage{array}
\usepackage{graphicx}
\usepackage{tikz}
\usepackage{hyperref}
\usepackage{cleveref}
\usepackage{fancybox}

\hypersetup{colorlinks=true,linkcolor=blue,urlcolor=blue,citecolor=blue}

% ---- Partial-derivative shorthand ----
\newcommand{\pds}[1]{\partial_{\sigma_{#1}}}   % d/d(sigma_n)

% ---- Vector/tensor and misc convenience ----
\newcommand{\vect}[1]{\mathbf{#1}}

% ---- Code-facing typography ----
\newcommand{\code}[1]{\texttt{#1}}
\newcommand{\jacparam}[1]{\code{#1}}

\title{\textbf{Analytical Resistivity Jacobian\\
       for a 1D Layered VTI Electromagnetic Medium\\}
       \large Public Version}
\author{Wouter Deleersnyder}
\date{\today}


\setlength{\parindent}{0pt}

\begin{document}
%\maketitle

\thispagestyle{empty}

\vspace*{3cm}

{\centering
	\fbox{
		\begin{minipage}{0.9\textwidth}
			\centering
			\vspace{0.5cm}
			{\huge \textbf{Analytical Resistivity Jacobian\\
					for a 1D Layered VTI Electromagnetic Medium}}
			\vspace{0.3cm}
			
			%{\large Subtitle goes here}
			\vspace{0.5cm}
		\end{minipage}
	}
	
	\vspace{2cm}
	
	{\large Wouter Deleersnyder}
	
	\vspace{0.5cm}
	
	Public Version
	
	\vspace{0.5cm}
	
	{\normalsize \today}
	
	\par}


	\vspace{5.5cm}
This document presents the full analytical derivation of the Jacobian for resistivity. Derivations for the remaining model parameters are intentionally not included here and are not planned for public release. This material is shared only upon individual request, in order to retain some control over how and by whom it is used.\\

If you are interested in the derivations for the other parameters, or in discussing potential applications, extensions, or collaboration, please feel free to reach out to me at\\ \url{www.wouterdeleersnyder.be}


\vspace*{\fill}

\newpage

\tableofcontents

\chapter{Introduction}
\section{Scope and how to read this document}


The electromagnetic field radiated by a dipole source in horizontally layered, homogeneous media can be expressed in the wavenumber domain, following a Hankel transform, in terms of up- and down-going TM- and TE-mode reflection coefficients and a small set of source-layer propagation factors. Differentiating this closed-form expression with respect to the conductivity of a given layer yields the conductivity sensitivity analytically, eliminating the need for finite-difference perturbations of the forward model.\\

\Cref{sec:model} establishes the notation and presents the constitutive relations linking conductivity to the two complex per-layer quantities—admittivity and impedivity—used by the forward model. \Cref{sec:tmte} introduces the reflection-coefficient and propagator framework before differentiation, which is then developed analytically in \Cref{sec:res}. The main result, presented in \Cref{sec:res}, is the resistivity Jacobian. \Cref{app:notation} provides a consolidated reference for the notation and symbols used throughout.\\

\emph{Sign convention:} \(e^{+i\omega t}\) is used throughout, so that \(s := i\omega\) with \(\omega = 2\pi f\) the angular frequency.

% ==============================================================
\section{Layered VTI model and constitutive relations}
\label{sec:model}
% ==============================================================

\subsection{Geometry}

The subsurface is represented by \(N\) horizontal, homogeneous VTI
(vertical transverse isotropic) layers separated by planar interfaces at
\(z_1 < z_2 < \dots < z_{N-1}\). The upper and lower half-spaces are
defined by \(z_0=-\infty\) and \(z_N=+\infty\), respectively. Layer \(n\)
extends over \(z_{n-1}<z<z_n\) and, for an interior layer, has thickness
\(d_n=z_n-z_{n-1}\). The source is located at depth
\(z_0^{\mathrm{src}}\) in layer \(s\), while the receiver is positioned at
depth \(z\). 

\subsection{Constitutive parameters}
\label{sec:constitutive}

Each layer \(n\) is described by real, frequency-independent material
properties: the horizontal and vertical resistivities
\(\rho_n,\rho_n^v\) (equivalently the conductivities), the horizontal and
vertical relative electric permittivities
\(\varepsilon_n,\varepsilon_n^v\), and the horizontal and vertical relative
magnetic permeabilities \(\mu_n,\mu_n^v\). From these quantities, four
complex frequency-dependent \emph{admittivity}/\emph{impedivity}
parameters are defined, with \(s:=i\omega\):
\begin{align}
    \eta_n   &= \sigma_n + s\,\varepsilon_n
             = \frac{1}{\rho_n} + i\omega\varepsilon_0\,\varepsilon_n,
             &\eta^v_n  &= \sigma^v_n + s\,\varepsilon^v_n
             = \frac{1}{\rho_n\,\lambda_n^2} + i\omega\varepsilon_0\,\varepsilon^v_n,
    \label{eq:eta-def} \\
    \zeta_n  &= s\mu_n = i\omega\mu_0\,\mu_n,
             &\zeta^v_n &= s\mu^v_n = i\omega\mu_0\,\mu^v_n,
    \label{eq:zeta-def}
\end{align}
Here \(\lambda_n=\sigma_n/\sigma_n^v=\sqrt{\rho_n^v/\rho_n}\geq1\)
denotes the dimensionless \emph{anisotropy factor};
\(\varepsilon_0\) and \(\mu_0\) are the vacuum permittivity and
permeability. The quantities \(\eta_n,\eta_n^v\) are referred to as
\emph{admittivities} and govern the electric response, whereas
\(\zeta_n,\zeta_n^v\) are \emph{impedivities} associated with the magnetic
response. Together they form the \emph{eta/zeta seeds} from which the
forward model and its derivatives are constructed. For the conductivity
derivatives in \Cref{sec:res}, only
\(\partial\eta_n/\partial\sigma_n=1\) and
\(\partial\eta_n^v/\partial\sigma_n^v=1\) are required, with the remaining
constitutive parameters kept fixed.

\subsection{TM/TE propagation constants}
\label{sec:propagation-constants}

For a planar layered medium, Maxwell's equations can be separated into
transverse-magnetic (TM) and transverse-electric (TE) modes. The TM mode
is associated with the vertical electric field, whereas the TE mode is
associated with the vertical magnetic field. Using the horizontal
wavenumber \(\kappa\) as the Hankel-transform integration variable, define
\begin{align}
    \gamma_n  &= \sqrt{\zeta_n\,\eta^v_n\vphantom{)}}\,, &
    \bar\gamma_n &= \sqrt{\eta_n\,\zeta^v_n\vphantom{)}}\,,
    \label{eq:gamma-def}\\
    \Gamma_n &= \left(\frac{\eta_n}{\eta^v_n}\right)^{1/2}
                \left(\kappa^2+\gamma_n^2\right)^{1/2},
             &
    \bar\Gamma_n &= \left(\frac{\zeta_n}{\zeta^v_n}\right)^{1/2}
                \left(\kappa^2+\bar\gamma_n^2\right)^{1/2}.
    \label{eq:Gamma-def}
\end{align}
\(\Gamma_n\) and \(\bar\Gamma_n\) are the TM and TE propagation
constants, respectively. In the isotropic limit,
\(\eta_n=\eta_n^v\) and \(\zeta_n=\zeta_n^v\), both reduce to the standard
diffusion expression \(\sqrt{\kappa^2+s\mu_n\eta_n}\). The two modes,
however, have different dependencies on the electric properties.
\(\bar\gamma_n\) contains the horizontal quantity \(\eta_n\) but not
\(\eta_n^v\), whereas \(\Gamma_n\) depends on both through \(\gamma_n\)
and the factor \((\eta_n/\eta_n^v)^{1/2}\). Consequently, the TE mode is
independent of vertical conductivity,
\(\partial\bar\Gamma_n/\partial\sigma_n^v\equiv0\), as shown in
\Cref{sec:res}, \eqref{eq:1.21}, while the TM mode remains sensitive to
both horizontal and vertical properties.

% ==============================================================
\section{TM/TE decomposition, reflection coefficients, propagators}
\label{sec:tmte}
% ==============================================================

This section summarizes the \emph{primal}, non-Jacobian quantities
that are differentiated in \Cref{sec:res}. These consist of the local and
global reflection coefficients, the source-layer field propagators, and
the TM/TE amplitude factors used to assemble the wavenumber-domain
Green's functions.

\subsection{Local reflection coefficients}
\label{sec:local-refl-ch3}

At the interface separating layers \(n\) and \(n+1\), the corresponding
\emph{local} single-interface reflection coefficients for the TE and TM
modes are
\begin{equation}
    r_n^{TE} = \frac{\bar\Gamma_n - \bar\Gamma_{n+1}}{\bar\Gamma_n + \bar\Gamma_{n+1}},
    \qquad
    r_n^{TM} = \frac{\eta_{n+1}\Gamma_n - \eta_n\Gamma_{n+1}}
                    {\eta_{n+1}\Gamma_n + \eta_n\Gamma_{n+1}}.
    \label{eq:ch3-local-refl}
\end{equation}
These expressions follow from enforcing continuity of the tangential
fields across the interface. The TE coefficient has the usual symmetric
mismatch form, while the TM coefficient contains the additional
\(\eta_{n+1}/\eta_n\) weighting because its characteristic impedance
depends on \(\Gamma_n/\eta_n\) rather than on \(\Gamma_n\) alone.

\subsection{Global reflection coefficients}
\label{sec:global-refl-ch3}

Multiple reflections within the layered medium are incorporated using
the standard recursive Bremmer-series formulation. The global reflection
coefficient \(R_n^{TE,TM}\), evaluated from the top of layer \(n+1\), is
then given by
\begin{equation}
    R_n^{TE,TM} =
        \frac{r_n^{TE,TM} + R_{n+1}^{TE,TM}\,E_{n+1}}
             {1 + r_n^{TE,TM}\,R_{n+1}^{TE,TM}\,E_{n+1}},
    \qquad
    E_{n+1} = \exp(-2\Gamma_{n+1}^{(\cdot)}\,d_{n+1}),
    \label{eq:ch3-global-refl}
\end{equation}
Here \(\Gamma^{(\cdot)}\) denotes \(\Gamma\) for TM and
\(\bar\Gamma\) for TE. The recursion starts from the bottom half-space
with \(R_N=0\) and proceeds upward; the analogous construction is used for
the opposite propagation direction. Since the contribution of each layer
contains the exponentially decaying factor \(E_{n+1}\), deeper interfaces
are increasingly attenuated. This is the physical basis for treating a
conductivity perturbation as a local modification that is propagated
outward through a sequence of
\(\partial R_k/\partial R_{k+1}\) factors
(\Cref{sec:res}, \eqref{eq:1.27}, \eqref{eq:1.33}).

\subsection{Source-layer field propagators and amplitude factors}
\label{sec:propagators-ch3}

Inside source layer \(s\), the vertical field dependence is represented
by two exponential propagators referenced to the receiver depth \(z\) and
the bounding interfaces \(z_{s-1}\) and \(z_s\):
\begin{equation}
    W_s^u = \exp\!\left[-\Gamma_s(z_s - z)\right], \qquad
    W_s^d = \exp\!\left[-\Gamma_s(z - z_{s-1})\right],
    \label{eq:ch3-propagators}
\end{equation}
The TE propagators \(\bar W_s^{u,d}\) have the same form after replacing
\(\Gamma_s\) by \(\bar\Gamma_s\). The upward- and downward-going amplitude
factors combine the source-layer global reflection coefficients
\(R_s^-\) and \(R_s^+\) with the multiple-reflection denominator
\begin{equation}
    M_s = 1 - R_s^-\,R_s^+\,\exp(-2\Gamma_s\,d_s)
    \label{eq:ch3-Ms}
\end{equation}
are
\begin{align}
    P_s^{u\pm} &= \frac{R_s^+}{M_s}
        \left(\exp(-\Gamma_s\,d^+) \pm R_s^-\exp\!\left[-\Gamma_s(d_s+d^-)\right]\right),
    \label{eq:ch3-Pu}\\
    P_s^{d\pm} &= \frac{R_s^-}{M_s}
        \left(\exp(-\Gamma_s\,d^-) \pm R_s^+\exp\!\left[-\Gamma_s(d_s+d^+)\right]\right),
    \label{eq:ch3-Pd}
\end{align}
where \(d^+=z_s-z_0^{\mathrm{src}}\) and
\(d^-=z_0^{\mathrm{src}}-z_{s-1}\). The corresponding TE factors
\(\bar P_s^{u\pm},\bar P_s^{d\pm}\) are obtained using
\(R^\pm\to R^{TE,\pm}\) and \(\Gamma_s\to\bar\Gamma_s\). The TM and TE
integral kernels entering the wavenumber-domain Green's functions
(\Cref{sec:res}, \eqref{eq:1.41}, \eqref{eq:1.75}, \eqref{eq:1.84}) are
assembled from these building blocks:
\[
    \tilde g^{TM}_{zh;s} = P_s^{u-}W_s^u - P_s^{d-}W_s^d, \qquad
    \tilde g^{TM}_{hh;s} = P_s^{u+}W_s^u + P_s^{d+}W_s^d, \qquad
    \tilde g^{TE}_{zz;s} = \bar P_s^{u+}\bar W_s^u + \bar P_s^{d+}\bar W_s^d.
\]

% ==============================================================
\chapter{Gradients with respect to resistivity/conductivity}
\label{sec:res}
% ==============================================================

\textbf{Literature.} Streich et al.\ (2011), \emph{Sensitivity of
controlled-source electromagnetic fields in planarly layered media};
Patriarca et al.\ (2011), \emph{Gradient based technique for
electromagnetic layered earth model data inversion}.

\section{Introduction}
Conductivity is commonly treated as either isotropic or VTI-anisotropic,
in which case the horizontal and vertical conductivities are different
(Everett and Constable 1999; Newman et al.\ 2010). As with the EM field
components themselves, the corresponding derivatives can be written as
Bessel-function integrals. These integrals can be evaluated efficiently
with fast Hankel-transform digital filters
(e.g.\ Johansen and S{\o}rensen 1979; Anderson 1982; Christensen 1990),
so that the computational effort required for sensitivities remains
comparable to that of the forward-field calculation.

For a planar layered medium, Maxwell's equations likewise allow the EM
field to be separated into TE and TM contributions
(\Cref{sec:tmte}).

\section{EM field derivatives}

For an arbitrary field component \(F\), the sensitivity \(J_n^F\) is
defined as the derivative with respect to the conductivity of layer \(n\),
\begin{equation}
    J^F_n = \frac{\partial F}{\partial \sigma_n}
    \label{eq:1.1}
\end{equation}
Alternatively, derivatives may be taken with respect to \(\ln\sigma\),
which is useful for imposing bounds while guaranteeing positive
conductivity values. The derivations below use the constitutive relations
introduced in \Cref{sec:model}:
\begin{align}
    s     &= i\omega                                                              \label{eq:1.2}  \\
    \zeta &= s\mu                                                                 \label{eq:1.3}  \\
    \zeta^v &= s\mu^v                                                             \label{eq:1.4}  \\
    \eta  &= \sigma + s\varepsilon                                                \label{eq:1.5}  \\
    \eta^v &= \sigma^v + s\varepsilon^v                                           \label{eq:1.6}  \\
    \gamma  &= \sqrt{\zeta\,\eta^v}                                               \label{eq:1.7}  \\
    \bar{\gamma} &= \sqrt{\eta\,\zeta^v}                                          \label{eq:1.8}  \\
    \Gamma  &= \left(\frac{\eta}{\eta^v}\right)^{1/2}\!\left(\kappa^2+\gamma^2\right)^{1/2}
                                                                                  \label{eq:1.9}  \\
    \bar{\Gamma} &= \left(\frac{\zeta}{\zeta^v}\right)^{1/2}\!\left(\kappa^2+\bar{\gamma}^2\right)^{1/2}
                                                                                  \label{eq:1.10}
\end{align}
We now derive explicit analytical expressions for the EM-field
sensitivities in a one-dimensional isotropic or VTI layered medium.

\subsection{Derivatives of the local reflection coefficients}
\label{sec:local_refl}

The TE- and TM-mode local reflection coefficients are given by
\begin{equation}
    r_n^{TE} = \frac{\Gamma_n - \Gamma_{n+1}}{\Gamma_n + \Gamma_{n+1}}
    \label{eq:1.11}
\end{equation}
and
\begin{equation}
    r_n^{TM} = \frac{\eta_{n+1}\,\Gamma_n - \eta_n\,\Gamma_{n+1}}
                    {\eta_{n+1}\,\Gamma_n + \eta_n\,\Gamma_{n+1}}
    \label{eq:1.12}
\end{equation}
Differentiating these expressions gives
\begin{align}
    \frac{\partial r_n^{TE}}{\partial \sigma_n}
    &= \pds{n}\!\left(\frac{\Gamma_n - \Gamma_{n+1}}{\Gamma_n + \Gamma_{n+1}}\right)
     = \frac{\pds{n}\Gamma_n}{\Gamma_n + \Gamma_{n+1}}
       + \frac{-(\Gamma_n - \Gamma_{n+1})\,(\pds{n}\Gamma_n)}{(\Gamma_n + \Gamma_{n+1})^2}
     = \frac{(\pds{n}\Gamma_n)(1 - r_n^{TE})}{\Gamma_n + \Gamma_{n+1}}
    \label{eq:1.13}
\end{align}
and
\begin{align}
    \frac{\partial r_n^{TM}}{\partial \sigma_n}
    &= \frac{\partial r_n^{TM}}{\partial \Gamma_n}\frac{\partial \Gamma_n}{\partial \sigma_n}
     + \frac{\partial r_n^{TM}}{\partial \eta_n}\frac{\partial \eta_n}{\partial \sigma_n}
     = \pds{n}\!\left(\frac{\eta_{n+1}\Gamma_n - \eta_n\Gamma_{n+1}}
                          {\eta_{n+1}\Gamma_n + \eta_n\Gamma_{n+1}}\right) \nonumber \\
    &= \frac{\eta_{n+1}\,\pds{n}\Gamma_n - \Gamma_{n+1}
             - (\eta_{n+1}\Gamma_n - \eta_n\Gamma_{n+1})(\eta_{n+1}\,\pds{n}\Gamma_n + \Gamma_{n+1})}
            {(\eta_{n+1}\Gamma_n + \eta_n\Gamma_{n+1})^2} \nonumber \\
    &= \frac{\eta_{n+1}(\pds{n}\Gamma_n)(1 - r_n^{TM}) - \Gamma_{n+1}(1 + r_n^{TM})}
            {\eta_{n+1}\Gamma_n + \eta_n\Gamma_{n+1}}
    \label{eq:1.14}
\end{align}
Because the properties of layer \(n\) also enter the reflection
coefficient at boundary \(n-1\), the corresponding derivatives are also
required:
\begin{align}
    \frac{\partial r_{n-1}^{TE}}{\partial \sigma_n}
    &= \pds{n}\!\left(\frac{\Gamma_{n-1} - \Gamma_n}{\Gamma_{n-1} + \Gamma_n}\right)
     = \frac{-\pds{n}\Gamma_n}{\Gamma_{n-1} + \Gamma_n}
       + \frac{-(\Gamma_{n-1} - \Gamma_n)\,(\pds{n}\Gamma_n)}{(\Gamma_{n-1} + \Gamma_n)^2}
     = \frac{-(\pds{n}\Gamma_n)(1 + r_{n-1}^{TE})}{\Gamma_{n-1} + \Gamma_n}
    \label{eq:1.15}
\end{align}
and
\begin{align}
    \frac{\partial r_{n-1}^{TM}}{\partial \sigma_n}
    &= \pds{n}\!\left(\frac{\eta_n\Gamma_{n-1} - \eta_{n-1}\Gamma_n}
                          {\eta_n\Gamma_{n-1} + \eta_{n-1}\Gamma_n}\right) \nonumber \\
    &= \frac{\Gamma_{n-1} - \eta_{n-1}(\pds{n}\Gamma_n)
             - r_{n-1}^{TM}\!\left(\Gamma_{n-1} + \eta_{n-1}(\pds{n}\Gamma_n)\right)}
            {\eta_n\Gamma_{n-1} + \eta_{n-1}\Gamma_n} \nonumber \\
    &= \frac{-\eta_{n-1}(\pds{n}\Gamma_n)(1 + r_{n-1}^{TM}) + \Gamma_{n-1}(1 - r_{n-1}^{TM})}
            {\eta_n\Gamma_{n-1} + \eta_{n-1}\Gamma_n}
    \label{eq:1.16}
\end{align}
where \(\dfrac{\partial \eta_n}{\partial \sigma_n} = 1\) and
\begin{align}
    \frac{\partial \Gamma_n}{\partial \sigma_n}
    &= \pds{n}\!\left((\kappa^2 + \gamma^2)^{1/2}\right)
     = \pds{n}\!\left((\kappa^2 + s\mu_n\eta_n)^{1/2}\right)
     = \pds{n}\!\left((\kappa^2 + s\mu_n(\sigma_n + s\varepsilon_n))^{1/2}\right) \nonumber \\
    &= \pds{n}\!\left((\kappa^2 + s\mu_n\sigma_n + s^2\mu_n\varepsilon_n)^{1/2}\right)
     = \frac{s\mu_n}{2\Gamma_n}
    \label{eq:1.17}
\end{align}
for the isotropic case, \(\eta_n=\eta_n^v\). In the VTI formulation,
the corresponding derivatives of \(\Gamma\) and \(\bar\Gamma\) are
\begin{align}
    \frac{\partial \Gamma_n}{\partial \sigma_n}
    &= \pds{n}\!\left(\left(\frac{\eta_n}{\eta^v_n}\right)^{1/2}(\kappa^2 + \gamma^2)^{1/2}\right)
     = \frac{\Gamma_n}{2\eta_n}
    \label{eq:1.18} \\
    \frac{\partial \bar{\Gamma}_n}{\partial \sigma_n}
    &= \pds{n}\!\left(\left(\frac{\zeta_n}{\zeta^v_n}\right)^{1/2}(\kappa^2 + \bar{\gamma}^2)^{1/2}\right)
     = \left(\frac{\zeta^v_n}{\zeta_n}\right)^{1/2}\frac{s\mu^v_n}{2\bar{\Gamma}_n}
    \label{eq:1.19} \\
    \frac{\partial \Gamma_n}{\partial \sigma^v_n}
    &= \left(\frac{\eta^v_n}{\eta_n}\right)^{1/2}\frac{s\mu^v_n}{2\bar{\Gamma}_n}
       - \frac{\Gamma_n}{2\eta^v_n}
    \label{eq:1.20} \\
    \frac{\partial \bar{\Gamma}_n}{\partial \sigma^v_n}
    &= 0
    \label{eq:1.21}
\end{align}
The final relation implies
\(\partial r^{TE}/\partial\sigma_n^v=0\). Thus the TE-mode expressions
depend only on the horizontal medium properties, confirming the structural
asymmetry identified in \Cref{sec:propagation-constants}.

\subsection{Derivatives of the global reflection coefficients}
\label{sec:global_refl}

The upward- and downward-looking global reflection coefficients
\(R_n^{TE,TM\pm}\) have the same general recursive structure for the TE
and TM modes:
\begin{equation}
    R_n^{TE,TM} =
        \frac{r_n^{TE,TM} + R_{n+1}^{TE,TM}\exp(-2\Gamma_{n+1}\,d_{n+1})}
             {1 + r_n^{TE,TM}\,R_{n+1}^{TE,TM}\exp(-2\Gamma_{n+1}\,d_{n+1})}
    \label{eq:1.22}
\end{equation}
For a perturbation of layer \(n\), the direct dependence enters through
\(r_n\), allowing \(\partial R_n/\partial\sigma_n\) to be written as
\begin{equation}
    \frac{\partial R_n}{\partial \sigma_n} =
        \frac{\partial R_n}{\partial r_n}\,\frac{\partial r_n}{\partial \sigma_n}
    \label{eq:1.23}
\end{equation}
where the required partial derivative with respect to \(r_n\) is
\begin{equation}
    \frac{\partial R_n}{\partial r_n} =
        \frac{1 - R_n\,R_{n+1}\exp(-2\Gamma_{n+1}\,d_{n+1})}
             {1 + r_n\,R_{n+1}\exp(-2\Gamma_{n+1}\,d_{n+1})}
    \label{eq:1.24}
\end{equation}
When the recursion is continued to layer \(n-1\), the derivative
\(\partial R_{n-1}/\partial\sigma_n\) must be evaluated:
\begin{equation}
    R_{n-1} =
        \frac{r_{n-1} + R_n\exp(-2\Gamma_n\,d_n)}
             {1 + r_{n-1}\,R_n\exp(-2\Gamma_n\,d_n)}
    \label{eq:1.25}
\end{equation}
At this stage, \(r_{n-1}\), \(R_n\), and \(E_n^2\) all carry a
dependence on \(\sigma_n\):
\begin{equation}
    \frac{\partial R_{n-1}}{\partial \sigma_n} =
        \frac{\partial R_{n-1}}{\partial r_{n-1}}\,\frac{\partial r_{n-1}}{\partial \sigma_n}
      + \frac{\partial R_{n-1}}{\partial R_n}\,\frac{\partial R_n}{\partial \sigma_n}
      + \frac{\partial R_{n-1}}{\partial E_n^2}\,\frac{\partial E_n^2}{\partial \sigma_n}
    \label{eq:1.26}
\end{equation}
with \(E_n^2 = \exp(-2\Gamma_n\,d_n)\) and with
\begin{align}
    \frac{\partial R_{n-1}}{\partial r_{n-1}}
    &= \frac{1 - R_{n-1}\,R_n\exp(-2\Gamma_n\,d_n)}
            {1 + r_{n-1}\,R_n\exp(-2\Gamma_n\,d_n)}
    \label{eq:1.27} \\
    \frac{\partial R_{n-1}}{\partial R_n}
    &= \frac{(1 - r_{n-1}\,R_{n-1})\exp(-2\Gamma_n\,d_n)}
            {1 + r_{n-1}\,R_n\exp(-2\Gamma_n\,d_n)}
    \label{eq:1.28} \\
    \frac{\partial R_{n-1}}{\partial E_n^2}
    &= \frac{R_n(1 - r_{n-1}\,R_{n-1})}
            {1 + r_{n-1}\,R_n\exp(-2\Gamma_n\,d_n)}
    \label{eq:1.29} \\
    \frac{\partial E_n^2}{\partial \sigma_n}
    &= -2d_n\exp(-2\Gamma_n\,d_n)\,(\pds{n}\Gamma_n)
    \label{eq:1.30}
\end{align}
\(\pds{n}\Gamma_n\) is given by Equations~\eqref{eq:1.18}--\eqref{eq:1.21}.
For the downward reflection \(\dfrac{\partial R_n}{\partial
\sigma_{n+1}} = 0\).

Continuing the recursion one layer farther gives the derivative
\(\partial R_{n-2}/\partial\sigma_n\):
\begin{equation}
    \frac{\partial R_{n-2}}{\partial \sigma_n} =
        \frac{\partial R_{n-2}}{\partial R_{n-1}}\,\frac{\partial R_{n-1}}{\partial \sigma_n}
    \label{eq:1.31}
\end{equation}
with
\begin{equation}
    \frac{\partial R_{n-2}}{\partial R_{n-1}} =
        \frac{(1 - r_{n-2}\,R_{n-2})\exp(-2\Gamma_{n-1}\,d_{n-1})}
             {1 + r_{n-2}\,R_{n-1}\exp(-2\Gamma_{n-1}\,d_{n-1})}
    \label{eq:1.32}
\end{equation}
Equation~\eqref{eq:1.31} shows that, for \(k\leq n-2\), the effect of
the conductivity perturbation is propagated recursively through the global
reflection coefficients, giving
\begin{equation}
    \frac{\partial R_k}{\partial \sigma_n} =
        \frac{\partial R_k}{\partial R_{k+1}}\,\frac{\partial R_{k+1}}{\partial \sigma_n}
    \label{eq:1.33}
\end{equation}
The general pattern is therefore to introduce the perturbation in the
local reflection coefficient associated with layer \(n\), and then
propagate its effect through the stack one layer at a time via
\(\partial R_k/\partial R_{k+1}\). The same structure applies to other
medium-parameter Jacobians, with an appropriately modified driving term;
here it is specialized to conductivity.

\subsection{Derivatives of the EM field components}
\label{sec:field_comps}

The vertical electric field component is
\begin{equation}
    \hat{G}^{ze;xe;s}(\mathbf{x},\mathbf{x}_0,\omega)
    = \hat{G}^{ze;xe;s:i}(\mathbf{x}-\mathbf{x}_0,\omega)
    + \frac{\cos\phi}{4\pi\eta^v_s}
      \int_{0}^{\infty} \tilde{g}^{TM}_{zh;s}\,J_1(\kappa r)\,\kappa^2\,\mathrm{d}\kappa
    \label{eq:1.34}
\end{equation}
and its derivative follows as:
\begin{align}
    \frac{\partial \hat{G}^{ze;xe;s}(\mathbf{x},\mathbf{x}_0,\omega)}{\partial \sigma_n}
    &= \frac{\partial}{\partial \sigma_n}\hat{G}^{ze;xe;s:i}(\mathbf{x}-\mathbf{x}_0,\omega)
     + \frac{\cos\phi}{4\pi\eta^v_s}
       \int_{0}^{\infty}
         \frac{\partial \tilde{g}^{TM}_{zh;s}}{\partial \sigma_n}
         J_1(\kappa r)\,\kappa^2\,\mathrm{d}\kappa
    \label{eq:1.35}
\end{align}
The remaining electric-field components are
\begin{align}
    \hat{G}^{xe;xe;s}(\mathbf{x},\mathbf{x}_0,\omega)
    &= \hat{G}^{xe;xe;s:i}(\mathbf{x}-\mathbf{x}_0,\omega)
     + \frac{1}{8\pi}
       \int_{0}^{\infty}\!
         \left(\frac{\Gamma_s\,\tilde{g}^{TM}_{hh;s}}{\eta_s}
              -\frac{\zeta_s\,\tilde{g}^{TE}_{zz;s}}{\bar{\Gamma}_s}\right)
         J_0(\kappa r)\,\kappa\,\mathrm{d}\kappa \nonumber \\
    &\quad
     - \frac{\cos(2\phi)}{8\pi}
       \int_{0}^{\infty}\!
         \left(\frac{\Gamma_s\,\tilde{g}^{TM}_{hh;s}}{\eta_s}
              +\frac{\zeta_s\,\tilde{g}^{TE}_{zz;s}}{\bar{\Gamma}_s}\right)
         J_2(\kappa r)\,\kappa\,\mathrm{d}\kappa
    \label{eq:1.36} \\
    \hat{G}^{ye;xe;s}(\mathbf{x},\mathbf{x}_0,\omega)
    &= \hat{G}^{ye;xe;s:i}(\mathbf{x}-\mathbf{x}_0,\omega)
     - \frac{\sin(2\phi)}{8\pi}
       \int_{0}^{\infty}\!
         \left(\frac{\Gamma_s\,\tilde{g}^{TM}_{hh;s}}{\eta_s}
              +\frac{\zeta_s\,\tilde{g}^{TE}_{zz;s}}{\bar{\Gamma}_s}\right)
         J_2(\kappa r)\,\kappa\,\mathrm{d}\kappa
    \label{eq:1.37}
\end{align}
Their derivatives are obtained in the same manner; the complete
expressions are given in \Cref{sec:horizontal_field_derivs}.

\begin{quote}
\small\itshape
\textbf{Worked example.} For the isotropic case (\(\sigma=\sigma^v\)) and
source/receiver both \(x\)-directed electric dipoles, Equation
\eqref{eq:1.36} is the field \(d = \hat G^{ee}_{xx;s}\) whose Jacobian
\(J=\partial d/\partial\sigma_n\) is a standard sanity-check case: since
Bessel functions and \(\kappa\) do not depend on \(\sigma_n\), the
derivative enters the integral directly and only the kernel derivatives
\(\partial\tilde g^{TM}_{hh;s}/\partial\sigma_n\),
\(\partial\tilde g^{TE}_{zz;s}/\partial\sigma_n\)
(\Cref{sec:hh_kernel}, \ref{sec:te_kernel}) need to be found, plus, when
\(n=s\), the direct dependence of the \(\Gamma_s/\eta_s\) and
\(\zeta_s/\bar\Gamma_s\) prefactors on \(\sigma_s\)
(\Cref{sec:horizontal_field_derivs}).
\end{quote}

\subsection{Derivatives of the integral coefficients}
\label{sec:integral_coeff}

The TE- and TM-mode integrals are given by
\begin{align}
    I^0_X &= \int_0^{\infty} \kappa\,J_0(\kappa r)\,g^X\,\mathrm{d}\kappa
    \label{eq:1.38} \\
    I^1_X &= \int_0^{\infty} J_1(\kappa r)\,g^X\,\mathrm{d}\kappa
    \label{eq:1.39}
\end{align}
where \(J_0(\kappa r)\) and \(J_1(\kappa r)\) are the zero- and
first-order Bessel functions and \(g^X\) are the integral kernels. The
derivatives of the Hankel integrals are given by
\begin{equation}
    \frac{\partial I^0_X}{\partial \sigma_n}
    = \frac{\partial}{\partial \sigma_n}\int_0^{\infty}\kappa\,J_0(\kappa r)\,g^X\,\mathrm{d}\kappa
    = \int_0^{\infty}\kappa\,J_0(\kappa r)\,\frac{\partial g^X}{\partial \sigma_n}\,\mathrm{d}\kappa
    \label{eq:1.40}
\end{equation}
To differentiate the integral coefficient
\(\tilde g^{TM}_{zh;s}\) in Equation~\eqref{eq:1.34}, the derivatives of
the field propagators and of the factors \(P_s^{u\pm}\) and
\(P_s^{d\pm}\) from \Cref{sec:tmte} are required:
\begin{equation}
    \tilde{g}^{TM}_{zh;s} = P_s^{u-}\,W_s^u - P_s^{d-}\,W_s^d
    \label{eq:1.41}
\end{equation}
The relevant propagators and amplitude factors are
\begin{align}
    W_n^u &= \exp\!\left[-\Gamma_n(z_n - z)\right] \label{eq:1.42} \\
    W_n^d &= \exp\!\left[-\Gamma_n(z - z_{n-1})\right] \label{eq:1.43}
\end{align}
and
\begin{align}
    P_s^{u\pm} &= \frac{R_s^+}{M_s}
        \left(\exp(-\Gamma_s\,d^+) \pm R_s^-\exp\!\left[-\Gamma_s(d_s+d^-)\right]\right)
    \label{eq:1.44} \\
    P_s^{d\pm} &= \frac{R_s^-}{M_s}
        \left(\exp(-\Gamma_s\,d^-) \pm R_s^+\exp\!\left[-\Gamma_s(d_s+d^+)\right]\right)
    \label{eq:1.45}
\end{align}
In these equations the factor \(M_s\) is given by
\begin{equation}
    M_s = 1 - R_s^-\,R_s^+\exp(-2\Gamma_s\,d_s)
    \label{eq:1.46}
\end{equation}
The derivative of the integral coefficient follows from the chain rule:
\begin{align}
    \frac{\partial \tilde{g}^{TM}_{zh;s}}{\partial \sigma_n}
    &= \frac{\partial \tilde{g}^{TM}_{zh;s}}{\partial P_s^{u-}}\,
       \frac{\partial P_s^{u-}}{\partial \sigma_n}
     + \frac{\partial \tilde{g}^{TM}_{zh;s}}{\partial W_s^u}\,
       \frac{\partial W_s^u}{\partial \sigma_n}
     + \frac{\partial \tilde{g}^{TM}_{zh;s}}{\partial P_s^{d-}}\,
       \frac{\partial P_s^{d-}}{\partial \sigma_n}
     + \frac{\partial \tilde{g}^{TM}_{zh;s}}{\partial W_s^d}\,
       \frac{\partial W_s^d}{\partial \sigma_n}
    \label{eq:1.47} \\
    &= W_s^u\,\frac{\partial P_s^{u-}}{\partial \sigma_n}
     + P_s^{u-}\,\frac{\partial W_s^u}{\partial \sigma_n}
     - W_s^d\,\frac{\partial P_s^{d-}}{\partial \sigma_n}
     - P_s^{d-}\,\frac{\partial W_s^d}{\partial \sigma_n}
    \label{eq:1.48}
\end{align}
where
\begin{align}
    \frac{\partial W_s^u}{\partial \sigma_n}
    &= -(z_n - z)\exp\!\left[-\Gamma_n(z_n-z)\right](\pds{n}\Gamma_n)
    \label{eq:1.49} \\
    \frac{\partial W_s^d}{\partial \sigma_n}
    &= -(z - z_{n-1})\exp\!\left[-\Gamma_n(z-z_{n-1})\right](\pds{n}\Gamma_n)
    \label{eq:1.50}
\end{align}
and
\begin{equation}
    \frac{\partial P_s^{u\pm}}{\partial \sigma_n}
    = \frac{\partial P_s^{u\pm}}{\partial R_s^+}\,\frac{\partial R_s^+}{\partial \sigma_n}
    + \frac{\partial P_s^{u\pm}}{\partial R_s^-}\,\frac{\partial R_s^-}{\partial \sigma_n}
    + \frac{\partial P_s^{u\pm}}{\partial M_s}\,\frac{\partial M_s}{\partial \sigma_n}
    + \frac{\partial P_s^{u\pm}}{\partial E_1^u}\,\frac{\partial E_1^u}{\partial \sigma_n}
    + \frac{\partial P_s^{u\pm}}{\partial E_2^u}\,\frac{\partial E_2^u}{\partial \sigma_n}
    \label{eq:1.51}
\end{equation}
The derivative \(\partial P_s^{u\pm}/\partial\sigma_n\) requires the
following structural partial derivatives:
\begin{align}
    \frac{\partial P_s^{u\pm}}{\partial R_s^+}
    &= \frac{1}{M_s}
       \left(\exp(-\Gamma_s\,d^+) \pm R_s^-\exp\!\left[-\Gamma_s(d_s+d^-)\right]\right)
    \label{eq:1.52} \\
    \frac{\partial P_s^{u\pm}}{\partial M_s}
    &= \frac{R_s^+}{M_s^2}
       \left(\exp(-\Gamma_s\,d^+) \pm R_s^-\exp\!\left[-\Gamma_s(d_s+d^-)\right]\right)
    \label{eq:1.53} \\
    \frac{\partial P_s^{u\pm}}{\partial R_s^-}
    &= \pm\frac{R_s^+}{M_s}\exp\!\left[-\Gamma_s(d_s+d^-)\right]
    \label{eq:1.54} \\
    \frac{\partial P_s^{u\pm}}{\partial E_1^u}
    &= \frac{R_s^+}{M_s}
    \label{eq:1.55} \\
    \frac{\partial P_s^{u\pm}}{\partial E_2^u}
    &= \pm\frac{R_s^-\,R_s^+}{M_s}
    \label{eq:1.56} \\
    \frac{\partial M_s}{\partial \sigma_n}
    &= \frac{\partial M_s}{\partial R_s^-}\,\frac{\partial R_s^-}{\partial \sigma_n}
     + \frac{\partial M_s}{\partial R_s^+}\,\frac{\partial R_s^+}{\partial \sigma_n}
     + \frac{\partial M_s}{\partial E_3}\,\frac{\partial E_3}{\partial \sigma_n}
    \label{eq:1.57} \\
    \frac{\partial M_s}{\partial R_s^\pm}
    &= R_s^\mp\exp(-2\Gamma_s\,d_s)
    \label{eq:1.58} \\
    \frac{\partial M_s}{\partial E_3}
    &= R_s^-\,R_s^+
    \label{eq:1.59} \\
    \frac{\partial E_3}{\partial \sigma_n}
    &= -2d_s\exp(-2\Gamma_s\,d_s)\,\pds{n}\Gamma_s
    \label{eq:1.60} \\
    \frac{\partial E_1^u}{\partial \sigma_n}
    &= -d^+\exp(-\Gamma_s\,d^+)\,\pds{n}\Gamma_s
    \label{eq:1.61} \\
    \frac{\partial E_2^u}{\partial \sigma_n}
    &= -(d_s+d^-)\exp\!\left[-\Gamma_s(d_s+d^-)\right]\pds{n}\Gamma_s
    \label{eq:1.62}
\end{align}
The quantities \(\partial R_s^\pm/\partial\sigma_n\) are obtained from
the reflection-coefficient recursion derived above.

The derivative of \(P_s^{d\pm}\) is obtained analogously:
\begin{equation}
    \frac{\partial P_s^{d\pm}}{\partial \sigma_n}
    = \frac{\partial P_s^{d\pm}}{\partial R_s^-}\,\frac{\partial R_s^-}{\partial \sigma_n}
    + \frac{\partial P_s^{d\pm}}{\partial R_s^+}\,\frac{\partial R_s^+}{\partial \sigma_n}
    + \frac{\partial P_s^{d\pm}}{\partial M_s}\,\frac{\partial M_s}{\partial \sigma_n}
    + \frac{\partial P_s^{d\pm}}{\partial E_1^d}\,\frac{\partial E_1^d}{\partial \sigma_n}
    + \frac{\partial P_s^{d\pm}}{\partial E_2^d}\,\frac{\partial E_2^d}{\partial \sigma_n}
    \label{eq:1.63}
\end{equation}
\begin{align}
    \frac{\partial P_s^{d\pm}}{\partial R_s^-}
    &= \frac{1}{M_s}
       \left(\exp(-\Gamma_s\,d^-) \pm R_s^+\exp\!\left[-\Gamma_s(d_s+d^+)\right]\right)
    \label{eq:1.64} \\
    \frac{\partial P_s^{d\pm}}{\partial M_s}
    &= \frac{R_s^-}{M_s^2}
       \left(\exp(-\Gamma_s\,d^-) \pm R_s^+\exp\!\left[-\Gamma_s(d_s+d^+)\right]\right)
    \label{eq:1.65} \\
    \frac{\partial P_s^{d\pm}}{\partial R_s^+}
    &= \pm\frac{R_s^-}{M_s}\exp\!\left[-\Gamma_s(d_s+d^+)\right]
    \label{eq:1.66} \\
    \frac{\partial P_s^{d\pm}}{\partial E_1^d}
    &= \frac{R_s^-}{M_s}
    \label{eq:1.67} \\
    \frac{\partial P_s^{d\pm}}{\partial E_2^d}
    &= \pm\frac{R_s^+\,R_s^-}{M_s}
    \label{eq:1.68} \\
    \frac{\partial M_s}{\partial \sigma_n}
    &= \frac{\partial M_s}{\partial R_s^+}\,\frac{\partial R_s^+}{\partial \sigma_n}
     + \frac{\partial M_s}{\partial R_s^-}\,\frac{\partial R_s^-}{\partial \sigma_n}
     + \frac{\partial M_s}{\partial E_3}\,\frac{\partial E_3}{\partial \sigma_n}
    \label{eq:1.69} \\
    \frac{\partial M_s}{\partial R_s^\pm}
    &= R_s^\mp\exp(-2\Gamma_s\,d_s)
    \label{eq:1.70} \\
    \frac{\partial M_s}{\partial E_3}
    &= R_s^-\,R_s^+
    \label{eq:1.71} \\
    \frac{\partial E_3}{\partial \sigma_n}
    &= -2d_s\exp(-2\Gamma_s\,d_s)\,\pds{n}\Gamma_s
    \label{eq:1.72} \\
    \frac{\partial E_1^d}{\partial \sigma_n}
    &= -d^-\exp(-\Gamma_s\,d^-)\,\pds{n}\Gamma_s
    \label{eq:1.73} \\
    \frac{\partial E_2^d}{\partial \sigma_n}
    &= -(d_s+d^+)\exp\!\left[-\Gamma_s(d_s+d^+)\right]\pds{n}\Gamma_s
    \label{eq:1.74}
\end{align}

\subsubsection{TM horizontal kernel \(\tilde{g}^{TM}_{hh;s}\) and its derivatives}
\label{sec:hh_kernel}

The horizontal TM kernel entering the \(x\)- and \(y\)-directed field
components in Equations~\eqref{eq:1.36}--\eqref{eq:1.37} is
\begin{equation}
    \tilde{g}^{TM}_{hh;s} = P_s^{u+}\,W_s^u + P_s^{d+}\,W_s^d,
    \label{eq:1.75}
\end{equation}
Here \(P_s^{u+}\) and \(P_s^{d+}\) correspond to the \(\pm=+\) cases of
Equations~\eqref{eq:1.44}--\eqref{eq:1.45}. This definition is consistent
with the source-depth identity
\(\partial\tilde{g}^{TM}_{zh;s}/\partial z_0=
\Gamma_s\tilde{g}^{TM}_{hh;s}\), which lies outside the present scope.
Applying the product rule gives
\begin{align}
    \frac{\partial\tilde{g}^{TM}_{hh;s}}{\partial\sigma_n}
    &= W_s^u\,\frac{\partial P_s^{u+}}{\partial\sigma_n}
     + P_s^{u+}\,\frac{\partial W_s^u}{\partial\sigma_n}
     + W_s^d\,\frac{\partial P_s^{d+}}{\partial\sigma_n}
     + P_s^{d+}\,\frac{\partial W_s^d}{\partial\sigma_n},
    \label{eq:1.76}
\end{align}
with \(\partial W_s^{u,d}/\partial\sigma_n\) as in
Equations~\eqref{eq:1.49}--\eqref{eq:1.50}. The chain rule for
\(\partial P_s^{u+}/\partial\sigma_n\) mirrors Equation~\eqref{eq:1.51}
with \(\pm=+\):
\begin{equation}
    \frac{\partial P_s^{u+}}{\partial\sigma_n}
    = \frac{\partial P_s^{u+}}{\partial R_s^+}\,\frac{\partial R_s^+}{\partial\sigma_n}
    + \frac{\partial P_s^{u+}}{\partial R_s^-}\,\frac{\partial R_s^-}{\partial\sigma_n}
    + \frac{\partial P_s^{u+}}{\partial M_s}\,\frac{\partial M_s}{\partial\sigma_n}
    + \frac{\partial P_s^{u+}}{\partial E_1^u}\,\frac{\partial E_1^u}{\partial\sigma_n}
    + \frac{\partial P_s^{u+}}{\partial E_2^u}\,\frac{\partial E_2^u}{\partial\sigma_n},
    \label{eq:1.77}
\end{equation}
with structural partial derivatives (Equations~\eqref{eq:1.52}--\eqref{eq:1.56}
evaluated at \(\pm=+\)):
\begin{align}
    \frac{\partial P_s^{u+}}{\partial R_s^+}
    &= \frac{E_1^u + R_s^-\,E_2^u}{M_s}, &
    \frac{\partial P_s^{u+}}{\partial M_s}
    &= \frac{R_s^+\!\left(E_1^u + R_s^-\,E_2^u\right)}{M_s^2}, \notag \\[2pt]
    \frac{\partial P_s^{u+}}{\partial R_s^-}
    &= \frac{R_s^+\,E_2^u}{M_s}, &
    \frac{\partial P_s^{u+}}{\partial E_1^u}
    &= \frac{R_s^+}{M_s}, &
    \frac{\partial P_s^{u+}}{\partial E_2^u}
    &= \frac{R_s^+\,R_s^-}{M_s}.
    \label{eq:1.78}
\end{align}
The derivative of \(P_s^{d+}\) is obtained in the same way from
Equation~\eqref{eq:1.63}, evaluated at \(\pm=+\):
\begin{equation}
    \frac{\partial P_s^{d+}}{\partial\sigma_n}
    = \frac{\partial P_s^{d+}}{\partial R_s^-}\,\frac{\partial R_s^-}{\partial\sigma_n}
    + \frac{\partial P_s^{d+}}{\partial R_s^+}\,\frac{\partial R_s^+}{\partial\sigma_n}
    + \frac{\partial P_s^{d+}}{\partial M_s}\,\frac{\partial M_s}{\partial\sigma_n}
    + \frac{\partial P_s^{d+}}{\partial E_1^d}\,\frac{\partial E_1^d}{\partial\sigma_n}
    + \frac{\partial P_s^{d+}}{\partial E_2^d}\,\frac{\partial E_2^d}{\partial\sigma_n},
    \label{eq:1.79}
\end{equation}
\begin{align}
    \frac{\partial P_s^{d+}}{\partial R_s^-}
    &= \frac{E_1^d + R_s^+\,E_2^d}{M_s}, &
    \frac{\partial P_s^{d+}}{\partial M_s}
    &= \frac{R_s^-\!\left(E_1^d + R_s^+\,E_2^d\right)}{M_s^2}, \notag \\[2pt]
    \frac{\partial P_s^{d+}}{\partial R_s^+}
    &= \frac{R_s^-\,E_2^d}{M_s}, &
    \frac{\partial P_s^{d+}}{\partial E_1^d}
    &= \frac{R_s^-}{M_s}, &
    \frac{\partial P_s^{d+}}{\partial E_2^d}
    &= \frac{R_s^+\,R_s^-}{M_s}.
    \label{eq:1.80}
\end{align}
The derivatives \(\partial M_s/\partial\sigma_n\) and
\(\partial E_i/\partial\sigma_n\) are identical to
Equations~\eqref{eq:1.57}--\eqref{eq:1.74}.

\subsubsection{TE kernel \(\tilde{g}^{TE}_{zz;s}\) and its derivatives}
\label{sec:te_kernel}

The TE kernel in Equations~\eqref{eq:1.36}--\eqref{eq:1.37} has the
same algebraic structure as the TM kernel, with \(\Gamma_s\) replaced by
\(\bar\Gamma_s\) and TE reflection coefficients used throughout. Define
the TE field propagators
\begin{align}
    \bar{W}_s^u &= \exp\!\left[-\bar\Gamma_s(z_s - z)\right], &
    \bar{W}_s^d &= \exp\!\left[-\bar\Gamma_s(z - z_{s-1})\right],
    \label{eq:1.81}
\end{align}
the TE multiple-reflection factor
\begin{equation}
    \bar{M}_s = 1 - R_s^{TE,-}\,R_s^{TE,+}\exp(-2\bar\Gamma_s\,d_s),
    \label{eq:1.82}
\end{equation}
the TE source-layer exponentials \(\bar{E}_1^u = e^{-\bar\Gamma_s d^+}\),
\(\bar{E}_2^u = e^{-\bar\Gamma_s(d_s+d^-)}\), \(\bar{E}_3 = e^{-2\bar\Gamma_s d_s}\),
\(\bar{E}_1^d = e^{-\bar\Gamma_s d^-}\), \(\bar{E}_2^d = e^{-\bar\Gamma_s(d_s+d^+)}\),
and the TE amplitude factors for the \(+\) variant:
\begin{align}
    \bar{P}_s^{u+}
    &= \frac{R_s^{TE,+}}{\bar{M}_s}
       \!\left(\bar{E}_1^u + R_s^{TE,-}\,\bar{E}_2^u\right), &
    \bar{P}_s^{d+}
    &= \frac{R_s^{TE,-}}{\bar{M}_s}
       \!\left(\bar{E}_1^d + R_s^{TE,+}\,\bar{E}_2^d\right).
    \label{eq:1.83}
\end{align}
The TE kernel can then be written as
\begin{equation}
    \tilde{g}^{TE}_{zz;s} = \bar{P}_s^{u+}\,\bar{W}_s^u + \bar{P}_s^{d+}\,\bar{W}_s^d.
    \label{eq:1.84}
\end{equation}
Because \(\bar\Gamma_n\) does not depend on \(\sigma_n^v\)
(Equation~\eqref{eq:1.21}), the complete TE kernel is independent of
vertical conductivity:
\begin{equation}
    \frac{\partial\tilde{g}^{TE}_{zz;s}}{\partial\sigma^v_n} = 0
    \qquad\text{for all }n.
    \label{eq:1.85}
\end{equation}
For horizontal conductivity, the product rule gives
\begin{align}
    \frac{\partial\tilde{g}^{TE}_{zz;s}}{\partial\sigma_{h,n}}
    &= \bar{W}_s^u\,\frac{\partial\bar{P}_s^{u+}}{\partial\sigma_{h,n}}
     + \bar{P}_s^{u+}\,\frac{\partial\bar{W}_s^u}{\partial\sigma_{h,n}}
     + \bar{W}_s^d\,\frac{\partial\bar{P}_s^{d+}}{\partial\sigma_{h,n}}
     + \bar{P}_s^{d+}\,\frac{\partial\bar{W}_s^d}{\partial\sigma_{h,n}},
    \label{eq:1.86}
\end{align}
where the propagator derivatives use \(\partial\bar\Gamma_n/\partial\sigma_{h,n}\)
from Equation~\eqref{eq:1.19}:
\begin{align}
    \frac{\partial\bar{W}_s^u}{\partial\sigma_{h,n}}
    &= -(z_s - z)\,\bar{W}_s^u\,\frac{\partial\bar\Gamma_n}{\partial\sigma_{h,n}}, &
    \frac{\partial\bar{W}_s^d}{\partial\sigma_{h,n}}
    &= -(z - z_{s-1})\,\bar{W}_s^d\,\frac{\partial\bar\Gamma_n}{\partial\sigma_{h,n}}.
    \label{eq:1.87}
\end{align}
The chain rule for \(\bar P_s^{u+}\) follows Equation~\eqref{eq:1.77}
after applying the substitutions \(R_s^\pm\to R_s^{TE,\pm}\),
\(M_s\to\bar M_s\), and \(E_i\to\bar E_i\):
\begin{align}
    \frac{\partial\bar{P}_s^{u+}}{\partial R_s^{TE,+}}
    &= \frac{\bar{E}_1^u + R_s^{TE,-}\bar{E}_2^u}{\bar{M}_s}, &
    \frac{\partial\bar{P}_s^{u+}}{\partial\bar{M}_s}
    &= \frac{R_s^{TE,+}\!\left(\bar{E}_1^u + R_s^{TE,-}\bar{E}_2^u\right)}{\bar{M}_s^2},
    \notag \\[2pt]
    \frac{\partial\bar{P}_s^{u+}}{\partial R_s^{TE,-}}
    &= \frac{R_s^{TE,+}\,\bar{E}_2^u}{\bar{M}_s}, &
    \frac{\partial\bar{P}_s^{u+}}{\partial\bar{E}_1^u}
    &= \frac{R_s^{TE,+}}{\bar{M}_s}, &
    \frac{\partial\bar{P}_s^{u+}}{\partial\bar{E}_2^u}
    &= \frac{R_s^{TE,+}\,R_s^{TE,-}}{\bar{M}_s},
    \label{eq:1.88}
\end{align}
\begin{align}
    \frac{\partial\bar{M}_s}{\partial\sigma_{h,n}}
    &= \frac{\partial\bar{M}_s}{\partial R_s^{TE,-}}\,\frac{\partial R_s^{TE,-}}{\partial\sigma_{h,n}}
     + \frac{\partial\bar{M}_s}{\partial R_s^{TE,+}}\,\frac{\partial R_s^{TE,+}}{\partial\sigma_{h,n}}
     + \frac{\partial\bar{M}_s}{\partial\bar{E}_3}\,\frac{\partial\bar{E}_3}{\partial\sigma_{h,n}},
    \label{eq:1.89}
\end{align}
\begin{align}
    \frac{\partial\bar{M}_s}{\partial R_s^{TE,\pm}}
    &= R_s^{TE,\mp}\,\bar{E}_3, &
    \frac{\partial\bar{M}_s}{\partial\bar{E}_3}
    &= R_s^{TE,-}\,R_s^{TE,+}, &
    \frac{\partial\bar{E}_3}{\partial\sigma_{h,n}}
    &= -2d_s\,\bar{E}_3\,\frac{\partial\bar\Gamma_s}{\partial\sigma_{h,n}},
    \label{eq:1.90}
\end{align}
\begin{align}
    \frac{\partial\bar{E}_1^u}{\partial\sigma_{h,n}}
    &= -d^+\,\bar{E}_1^u\,\frac{\partial\bar\Gamma_s}{\partial\sigma_{h,n}}, &
    \frac{\partial\bar{E}_2^u}{\partial\sigma_{h,n}}
    &= -(d_s+d^-)\,\bar{E}_2^u\,\frac{\partial\bar\Gamma_s}{\partial\sigma_{h,n}},
    \label{eq:1.91}
\end{align}
The corresponding expressions for \(\bar P_s^{d+}\) follow analogously
from Equation~\eqref{eq:1.79}:
\begin{align}
    \frac{\partial\bar{P}_s^{d+}}{\partial R_s^{TE,-}}
    &= \frac{\bar{E}_1^d + R_s^{TE,+}\bar{E}_2^d}{\bar{M}_s}, &
    \frac{\partial\bar{P}_s^{d+}}{\partial\bar{M}_s}
    &= \frac{R_s^{TE,-}\!\left(\bar{E}_1^d + R_s^{TE,+}\bar{E}_2^d\right)}{\bar{M}_s^2},
    \notag \\[2pt]
    \frac{\partial\bar{P}_s^{d+}}{\partial R_s^{TE,+}}
    &= \frac{R_s^{TE,-}\,\bar{E}_2^d}{\bar{M}_s}, &
    \frac{\partial\bar{P}_s^{d+}}{\partial\bar{E}_1^d}
    &= \frac{R_s^{TE,-}}{\bar{M}_s}, &
    \frac{\partial\bar{P}_s^{d+}}{\partial\bar{E}_2^d}
    &= \frac{R_s^{TE,+}\,R_s^{TE,-}}{\bar{M}_s},
    \label{eq:1.92}
\end{align}
\begin{align}
    \frac{\partial\bar{E}_1^d}{\partial\sigma_{h,n}}
    &= -d^-\,\bar{E}_1^d\,\frac{\partial\bar\Gamma_s}{\partial\sigma_{h,n}}, &
    \frac{\partial\bar{E}_2^d}{\partial\sigma_{h,n}}
    &= -(d_s+d^+)\,\bar{E}_2^d\,\frac{\partial\bar\Gamma_s}{\partial\sigma_{h,n}}.
    \label{eq:1.93}
\end{align}
The TE global-reflection derivatives
\(\partial R_s^{TE,\pm}/\partial\sigma_{h,n}\) follow the recursion in
\Cref{sec:global_refl}, with
\((\Gamma,r^{TM},R^{TM})\) replaced by
\((\bar\Gamma,r^{TE},R^{TE})\).

\subsubsection{Full derivatives of the horizontal field components}
\label{sec:horizontal_field_derivs}

The horizontal field components in
Equations~\eqref{eq:1.36}--\eqref{eq:1.37} contain products of the
source-layer prefactors \(\Gamma_s/\eta_s\) and
\(\zeta_s/\bar\Gamma_s\) with their respective kernels. If \(n=s\),
these prefactors also depend on \(\sigma_{\bullet,s}\) and therefore
contribute additional derivative terms. Applying the product rule gives:
\begin{equation}
    \frac{\partial}{\partial\sigma_{\bullet,n}}\!\left(
        \frac{\Gamma_s}{\eta_s}\,\tilde{g}^{TM}_{hh;s}\right)
    = \delta_{ns}\,\frac{\partial(\Gamma_s/\eta_s)}{\partial\sigma_{\bullet,s}}
      \,\tilde{g}^{TM}_{hh;s}
    + \frac{\Gamma_s}{\eta_s}\,
      \frac{\partial\tilde{g}^{TM}_{hh;s}}{\partial\sigma_{\bullet,n}},
    \label{eq:1.94}
\end{equation}
\begin{equation}
    \frac{\partial}{\partial\sigma_{\bullet,n}}\!\left(
        \frac{\zeta_s}{\bar\Gamma_s}\,\tilde{g}^{TE}_{zz;s}\right)
    = \delta_{ns}\,\frac{\partial(\zeta_s/\bar\Gamma_s)}{\partial\sigma_{\bullet,s}}
      \,\tilde{g}^{TE}_{zz;s}
    + \frac{\zeta_s}{\bar\Gamma_s}\,
      \frac{\partial\tilde{g}^{TE}_{zz;s}}{\partial\sigma_{\bullet,n}}.
    \label{eq:1.95}
\end{equation}
Because \(\zeta_s=s\mu_s\) has no conductivity dependence, the non-zero
prefactor derivatives reduce to the following expressions, using
Equations~\eqref{eq:1.18}--\eqref{eq:1.21}:
\begin{align}
    \frac{\partial(\Gamma_s/\eta_s)}{\partial\sigma_{h,s}}
    &= \frac{(\pds{s}\Gamma_s)\,\eta_s - \Gamma_s}{\eta_s^2}
     = \frac{\Gamma_s/2 - \Gamma_s}{\eta_s^2}
     = -\frac{\Gamma_s}{2\eta_s^2},
    \label{eq:1.96} \\[4pt]
    \frac{\partial(\Gamma_s/\eta_s)}{\partial\sigma^v_s}
    &= \frac{1}{\eta_s}\,\frac{\partial\Gamma_s}{\partial\sigma^v_s}
     = \frac{1}{\eta_s}
       \!\left[\left(\frac{\eta^v_s}{\eta_s}\right)^{\!1/2}
               \frac{s\mu^v_s}{2\bar\Gamma_s}
             - \frac{\Gamma_s}{2\eta^v_s}\right],
    \label{eq:1.97} \\[4pt]
    \frac{\partial(\zeta_s/\bar\Gamma_s)}{\partial\sigma_{h,s}}
    &= -\frac{\zeta_s}{\bar\Gamma_s^2}\,\frac{\partial\bar\Gamma_s}{\partial\sigma_{h,s}}
     = -\frac{\zeta_s}{\bar\Gamma_s^2}
       \left(\frac{\zeta^v_s}{\zeta_s}\right)^{\!1/2}\!\frac{s\mu^v_s}{2\bar\Gamma_s},
    \label{eq:1.98} \\[4pt]
    \frac{\partial(\zeta_s/\bar\Gamma_s)}{\partial\sigma^v_s}
    &= 0.
    \label{eq:1.99}
\end{align}
Combining the preceding terms yields the complete derivatives of the
horizontal field components:
\begin{align}
    \frac{\partial\hat{G}^{xe;xe;s}}{\partial\sigma_{\bullet,n}}
    &= \frac{\partial\hat{G}^{xe;xe;s:i}}{\partial\sigma_{\bullet,n}} \nonumber \\
    &\quad
    + \frac{1}{8\pi}\int_0^\infty\!\left[
        \delta_{ns}\frac{\partial(\Gamma_s/\eta_s)}{\partial\sigma_{\bullet,s}}
            \tilde{g}^{TM}_{hh;s}
      + \frac{\Gamma_s}{\eta_s}
            \frac{\partial\tilde{g}^{TM}_{hh;s}}{\partial\sigma_{\bullet,n}}
      - \delta_{ns}\frac{\partial(\zeta_s/\bar\Gamma_s)}{\partial\sigma_{\bullet,s}}
            \tilde{g}^{TE}_{zz;s}
      - \frac{\zeta_s}{\bar\Gamma_s}
            \frac{\partial\tilde{g}^{TE}_{zz;s}}{\partial\sigma_{\bullet,n}}
      \right]J_0(\kappa r)\,\kappa\,\mathrm{d}\kappa \nonumber \\
    &\quad
    - \frac{\cos(2\phi)}{8\pi}\int_0^\infty\!\left[
        \delta_{ns}\frac{\partial(\Gamma_s/\eta_s)}{\partial\sigma_{\bullet,s}}
            \tilde{g}^{TM}_{hh;s}
      + \frac{\Gamma_s}{\eta_s}
            \frac{\partial\tilde{g}^{TM}_{hh;s}}{\partial\sigma_{\bullet,n}}
      + \delta_{ns}\frac{\partial(\zeta_s/\bar\Gamma_s)}{\partial\sigma_{\bullet,s}}
            \tilde{g}^{TE}_{zz;s}
      + \frac{\zeta_s}{\bar\Gamma_s}
            \frac{\partial\tilde{g}^{TE}_{zz;s}}{\partial\sigma_{\bullet,n}}
      \right]J_2(\kappa r)\,\kappa\,\mathrm{d}\kappa,
    \label{eq:1.100}
\end{align}
\begin{align}
    \frac{\partial\hat{G}^{ye;xe;s}}{\partial\sigma_{\bullet,n}}
    &= \frac{\partial\hat{G}^{ye;xe;s:i}}{\partial\sigma_{\bullet,n}}
    - \frac{\sin(2\phi)}{8\pi}\int_0^\infty\!\left[
        \delta_{ns}\frac{\partial(\Gamma_s/\eta_s)}{\partial\sigma_{\bullet,s}}
            \tilde{g}^{TM}_{hh;s}
      + \frac{\Gamma_s}{\eta_s}
            \frac{\partial\tilde{g}^{TM}_{hh;s}}{\partial\sigma_{\bullet,n}}
      + \delta_{ns}\frac{\partial(\zeta_s/\bar\Gamma_s)}{\partial\sigma_{\bullet,s}}
            \tilde{g}^{TE}_{zz;s}
      + \frac{\zeta_s}{\bar\Gamma_s}
            \frac{\partial\tilde{g}^{TE}_{zz;s}}{\partial\sigma_{\bullet,n}}
      \right]J_2(\kappa r)\,\kappa\,\mathrm{d}\kappa.
    \label{eq:1.101}
\end{align}
For \(n\neq s\), the Kronecker-delta terms vanish and only the kernel
derivatives remain. When \(n=s\), the additional prefactor contributions
from \eqref{eq:1.96}--\eqref{eq:1.99} must also be included. The required
kernel derivatives are given in \Cref{sec:hh_kernel} and
\Cref{sec:te_kernel}.

\subsection{From \(\sigma\)-derivatives to a resistivity Jacobian}
\label{sec:res-code-crosscheck}

The derivatives above are expressed in terms of \(\sigma_n\) and
\(\sigma_n^v\), or equivalently \(\eta_n,\eta_n^v\) through
Equations~\eqref{eq:1.5}--\eqref{eq:1.6}. If the model is instead
parameterized by resistivity \(\rho_n=1/\sigma_n\) and anisotropy
\(\lambda_n=\sigma_n/\sigma_n^v\), the chain rule provides the conversion
from conductivity to resistivity derivatives:
\begin{equation}
    \eta_n = \frac{1}{\rho_n} + i\omega\varepsilon_0\varepsilon_n, \quad
    \eta^v_n = \frac{1}{\rho_n\lambda_n^2} + i\omega\varepsilon_0\varepsilon^v_n
    \quad\Longrightarrow\quad
    \frac{\partial\eta_n}{\partial\rho_n} = -\frac{1}{\rho_n^2}, \quad
    \frac{\partial\eta^v_n}{\partial\rho_n} = -\frac{1}{\rho_n^2\lambda_n^2}.
    \label{eq:res-code-seed}
\end{equation}
An implementation based directly on the eta/zeta seeds can therefore use
\(\partial\eta_n/\partial\rho_n\) and
\(\partial\eta_n^v/\partial\rho_n\), together with the conductivity
derivatives derived above, in the chain rule. In particular,
\(\partial F/\partial\rho_n\) is obtained by combining the horizontal and
vertical conductivity derivatives with
\(\partial\sigma_n/\partial\rho_n=-1/\rho_n^2\) and
\(\partial\sigma_n^v/\partial\rho_n=-1/(\rho_n^2\lambda_n^2)\) at fixed
\(\lambda_n\).

\chapter{Other gradients}

{\centering
	\fbox{
		\begin{minipage}{0.9\textwidth}
			\centering
			\vspace{0.5cm}
			
			If you are interested in the derivations for the other parameters, or in discussing potential applications, extensions, or collaboration, please feel free to reach out to me at \url{www.wouterdeleersnyder.be}.
			
			\vspace{0.5cm}
		\end{minipage}
	}
}

\section{Gradients with respect to the anisotropy factor (aniso)}

\section{Gradients with respect to electric permittivity (epermH, epermV)}

\section{Gradients with respect to magnetic permeability (mpermH, mpermV)}

\section{Gradients with respect to interface depths (depth, thickness)}

\section{Gradients with respect to source/receiver depth}

\section{Gradients with respect to offset (off)}

\section{Gradients with respect to frequency (freq)}

\section{Gradients with respect to horizontal source position}

\section{Gradients with respect to source/receiver orientation}

% ==============================================================
\appendix
\chapter{Notation}
\label{app:notation}
% ==============================================================

\section{Layer/medium parameters}

\begin{longtable}{lll}
\toprule
Symbol & Meaning & Units/notes \\
\midrule
\endhead
\(\rho_n\) & Horizontal resistivity & \(\Omega\)m; \(\rho_n=1/\sigma_n\) \\
\(\lambda_n=\sigma_n/\sigma^v_n\) & Anisotropy factor & \(=\sqrt{\rho^v_n/\rho_n}\geq 1\) \\
\(\varepsilon_n,\varepsilon^v_n\) & Relative electric permittivity & horizontal/vertical \\
\(\mu_n,\mu^v_n\) & Relative magnetic permeability & horizontal/vertical \\
\(\sigma_n=1/\rho_n,\ \sigma^v_n\) & Horizontal/vertical conductivity & S/m \\
\(\eta_n,\eta^v_n\) & Admittivity & \(=\sigma+i\omega\varepsilon_0\varepsilon\) \\
\(\zeta_n,\zeta^v_n\) & Impedivity & \(=i\omega\mu_0\mu\) \\
\(z_m\) & Interface depth & \\
\(d_m=z_m-z_{m-1}\) & Layer thickness & finite for interior layers \\
\(\omega=2\pi f\) & Angular frequency & \\
\(s=i\omega\) & Laplace/frequency variable & \(e^{+i\omega t}\) convention \\
\(\varepsilon_0,\mu_0\) & Vacuum permittivity/permeability & \\
\bottomrule
\end{longtable}

\section{Propagation, reflection and kernel quantities}

\begin{longtable}{ll}
\toprule
Symbol & Meaning \\
\midrule
\endhead
\(\kappa\) & Hankel-transform wavenumber \\
\(\Gamma_n\) & TM propagation constant (\Cref{sec:propagation-constants}) \\
\(\bar\Gamma_n\) & TE propagation constant \\
\(\gamma_n=\sqrt{\zeta_n\eta^v_n}\), \(\bar\gamma_n=\sqrt{\eta_n\zeta^v_n}\) & intermediate quantities \\
\(r_n^{TE},r_n^{TM}\) & Local reflection coefficients (\Cref{sec:tmte}) \\
\(R_n^{TE},R_n^{TM}\) & Global reflection coefficients \\
\(W_s^{u,d},\bar W_s^{u,d}\) & Source-layer field propagators \\
\(P_s^{u,d\pm},\bar P_s^{u,d\pm}\) & TM/TE amplitude factors \\
\(M_s,\bar M_s\) & Source-layer multiple factor \\
\(E_1^{u,d},E_2^{u,d},E_3\) & Source-layer exponentials \\
\(\tilde g^{TM}_{zh;s},\tilde g^{TM}_{hh;s},\tilde g^{TE}_{zz;s}\) & Integral (wavenumber-domain) kernels \\
\(\phi\) & Source-receiver azimuth \\
\(r\) & Source-receiver radial offset \\
\(z_0^{\mathrm{src}},z\) & Source/receiver depth \\
\(d^+,d^-\) & Source-to-interface distances \\
\bottomrule
\end{longtable}


\end{document}
